\title{Logical Complexes}
\author{
       Brian Droncheff
}

\documentclass[12pt]{report}
\usepackage{verbatim}
\usepackage{amsmath}
\usepackage{amssymb}
\usepackage{tikz-cd}
\usepackage{tikz}
\usepackage{tikz-3dplot}
\usepackage{scalerel}
\usepackage{stackengine}
\usepackage{mathrsfs}
\usepackage[all]{xy}
\usepackage{graphicx}
\usepackage{subcaption}
\usepackage{hyperref}
\usepackage{titlesec}     
\usepackage{changepage}
\usepackage{parallel,enumitem}
\usepackage[all]{xy}
\usepackage{floatrow}
\usepackage{glossaries}
\usepackage [autostyle, english = american]{csquotes}

\usetikzlibrary{calc,intersections}
\usetikzlibrary{positioning}
\usetikzlibrary{shapes.geometric,calc,through}

\newtheorem{theorem}{Theorem}[section]
\newtheorem{lemma}[theorem]{Lemma}
\newtheorem{proposition}[theorem]{Proposition}
\newtheorem{corollary}[theorem]{Corollary}
\newtheorem{example}[theorem]{Example}
\newtheorem{postulate}[theorem]{Postulate}
\newtheorem{diagram}[theorem]{Diagram}


\newenvironment{proof}[1][Proof]{\begin{trivlist}
\item[\hskip \labelsep {\bfseries #1}]}{\end{trivlist}}
\newenvironment{definition}[1][Definition]{\begin{trivlist}
\item[\hskip \labelsep {\bfseries #1}]}{\end{trivlist}}
\newenvironment{question}[1][Question]{\begin{trivlist}
\item[\hskip \labelsep {\bfseries #1}]}{\end{trivlist}}
\newenvironment{note}[1][Note]{\begin{trivlist}
\item[\hskip \labelsep {\bfseries #1}]}{\end{trivlist}}


\newcommand{\qed}{\nobreak \ifvmode \relax \else
      \ifdim\lastskip<1.5em \hskip-\lastskip
      \hskip1.5em plus0em minus0.5em \fi \nobreak
      \vrule height0.75em width0.5em depth0.25em\fi}
\newcommand{\noteend}{\nobreak \ifvmode \relax \else
      \ifdim\lastskip<1.5em \hskip-\lastskip
      \fi \nobreak
      \vrule height0.3em width0.3em depth0.25em\fi}
\newcommand{\defend}{$\spadesuit$}
\newcommand{\exend}{$\blacklozenge$}
\newcommand{\tb}{$\hspace*{.5cm}$}
\newcommand{\ltb}{$\hspace*{1cm}$}
\newcommand{\smsp}{$\hspace*{.1cm}$}
\newcommand{\smnl}{$\\[.25cm]$}
\newcommand{\dnl}{$\\[.5cm]$}
\newcommand{\Lim}[1]{\raisebox{0.5ex}{\scalebox{0.8}{$\displaystyle \lim_{#1}\;$}}}
\newcommand*{\inbk}[1]{$\langle {#1}|{#1} \rangle $}
\newcommand*{\otbk}[1]{$|{#1} \rangle \langle {#1}|$}
\newcommand*{\inbktwo}[2]{$\langle {#1}|{#2} \rangle $}
\newcommand*{\otbktwo}[2]{$|{#1} \rangle \langle {#2}|$}
\newcommand*{\otbkthree}[3]{${#1}|[#2]\rangle \langle [#2]|{#3} $}
\newcommand*{\inbkthree}[3]{$\langle {#1}|[#2]|{#3} \rangle $}
\newcommand{\rowtens}[2]{$|{#1} \rangle \otimes |{#2} \rangle$}
\newcommand{\coltens}[2]{$ \langle{#1}| \otimes  \langle{#2}|$}
\newcommand{\ket}[1]{$|{#1}\rangle$}
\newcommand{\bra}[1]{$\langle{#1}|$}
\newcommand{\tit}[1]{$\textit{#1}$}
\newcommand{\tbf}[1]{$\textbf{#1}$}
\newcommand*\circled[1]{\tikz[baseline=(char.base)]{
            \node[shape=circle,draw,inner sep=-.35pt] (char) {#1};}}
\newcommand{\ketrep}[1]{$[|{#1}\rangle]$}
\newcommand{\brarep}[1]{$[\langle{#1}|]$}
\newcommand{\ketreparg}[2]{${#1}[|\rangle]{#2}$}
\newcommand{\brareparg}[2]{${#1}[\langle|]{#2}$}
\newcommand{\comm}[2]{$ \left[{#1},{#2}\right]={#1}{#2}-{#2}{#1}$}
\newcommand{\tensormult}[4]{$\begin{bmatrix} #1#3 & #1#4 \\ #2#3 & #2#4 \end{bmatrix}$}
\newcommand{\tens}[2]{${#1} \otimes {#2}$}
          

%\def\subcap{\mathrel{\ThisStyle{\kern.7pt\stackinset{c}{-.3\LMpt}{b}{.8\LMpt}%
 % {\scalebox{1.5}[1]{$\SavedStyle\subset$}}%
 % {\scalebox{1}[1.5]{$\SavedStyle\cap$}}}\kern .3pt}}
\def\impax{\mathrel{\ThisStyle{\kern 0pt\stackinset{c}{0\LMpt}{b}{3\LMpt}%
  {\scalebox{1}[1]{$\SavedStyle\Leftrightarrow$}}%
  {\scalebox{1}[1]{$\SavedStyle\Updownarrow$}}}\kern 0pt}}
\makeatletter
\newcommand{\superimpose}[2]{%
  {\ooalign{$#1\@firstoftwo#2$\cr\hfil$#1\@secondoftwo#2$\hfil\cr}}}
\makeatother
\newcommand{\subsetintersection}{\mathpalette\superimpose{{\subset}{\cap}}}
\newenvironment{pindent}{\begin{adjustwidth}{1cm}{1cm}}{\end{adjustwidth}}

\MakeOuterQuote{"}
\begin{document}

Also note that for the logical metric defined on $\Updownarrow$, $(X \Updownarrow Y) \Updownarrow (Y \Updownarrow Z) = \partial (XYZ) \vee \partial^2(XYZ)$ or $\partial(XYZ) \Updownarrow \partial^2(XYZ)$. \smnl
\tb Consider the partitioned expression $D \vee G \vee B \vee C = (X \wedge Y \wedge \neg Z) \vee (X \wedge Y \wedge Z) \vee (Y \wedge \neg X \wedge \neg Z) \vee (Z \wedge \neg X \wedge \neg Y) = \mathcal{L}$. We will let the logical point for each $D,G,B,$ and $C$ be $\mathcal{V}^T(-)$ i.e. $\mathcal{V}^T(D)$ etc. as the logical points are those points which make the expression such as $D$ true. now for $\mathcal{V}^T(\mathcal{L})$ we have the points $\{(1,1,0),(1,1,1),(0,1,0),(0,0,1)\}$ corresponding to $\{\mathcal{V}^T(D),\mathcal{V}^T(G),\mathcal{V}^T(B),\mathcal{V}^T(C)\}$ which is a logical 3-simplex. Each of the logical points of the disjunctive terms of the form $(\mathcal{V}(X), \mathcal{V}(Y), \mathcal{V}(Z))$ in $\mathcal{L}$, for example the expression $(X \wedge Y \wedge \neg Z)$ has a logical point of $(\mathcal{V}(X), \mathcal{V}(Y), \mathcal{V}(Z)) = (1,1,0)$. The four possible 2-simplices are $R = \{(0,0,1),(1,1,1),(1,1,0)\}$, $S = \{(1,1,1),(1,1,0),(0,1,0)\}$, $T = \{(1,1,0),(0,1,0),(0,0,1)\}$, and $U = \{(0,0,1),(1,1,0),(0,1,0)\}$. Now as in simplical decomposition we verify the distance metric for elements of $R$:
\begin{center}
$((0,0,1) \Updownarrow (1,1,1)) \Updownarrow ((1,1,1) \Updownarrow (1,1,0)) = (1,1,1) = ((0,0,1) \Updownarrow (1,1,0))$ 
\end{center}
which is correct. \smnl
for $S$:
\begin{center}
$((1,1,1) \Updownarrow (1,1,0)) \Updownarrow ((1,1,0) \Updownarrow (0,1,0)) = (1,0,1) = ((1,1,1) \Updownarrow (0,1,0))$
\end{center}
likewise for $T$ and $U$. \smnl
\tb So we can create vectors for XOR of points (where XOR is the distance metric) and the XOR of two vectors give the third vector! What about planes then? \smnl
\tb Similar to the simplical decomposition in which $\partial(HIJ) = IJ - HJ + HI$ and $\partial^2(HIJ) = (J - I) - (J - H) + (I - H) = 0$, again we let the logial points be $\mathcal{V}^T(B) = (0,1,0)$, $\mathcal{V}^T(C) = (0,0,1)$, $\mathcal{V}^T(D) = (1,1,1)$, and $\mathcal{V}^T(G) = (1,1,1)$, just using $B,C,D,$ and $G$ for short; we have $\partial^2(CGD) = (GD) - (CD) + (CG) = ((1,1,0) \Updownarrow ((1,1,1))) \Updownarrow ((1,1,0) \Updownarrow (0,0,1)) \Updownarrow ((1,1,1) \Updownarrow (0,0,1)) = (0,0,1) \Updownarrow (1,1,1) \Updownarrow (1,1,0) = (0,0,0)$ \smnl
\tb Now note that for any points $R$ and $S$ that, logically if negation between points is $XOR$ as due to difference then as $\partial(RS) = S - R \equiv S \Updownarrow R$ then one interpretation of $-(RS)$ is $-(S - R) = R - S \equiv R \Updownarrow S$. Hence we can see that $RS = -(RS)$ hence we need not worry about negative signs. Another interpretation however might be that $-(RS) \equiv \neg(R \Updownarrow S) = R \Leftrightarrow S)$. We would then have that $\partial^2(CGD) = (GD) - (CD) + (CG) = (D \Updownarrow G) \Updownarrow (C \Leftrightarrow D) \Updownarrow (C \Updownarrow G) = ?$ However if we had that $-(C \Updownarrow D) = (D \Updownarrow C)$ for the interpretation that $-(CD) = (DC)$ then we have $\partial^2(CGD) = (D \Updownarrow G) \Updownarrow (D \Updownarrow C) \Updownarrow (C \Updownarrow G) = (D \Updownarrow D) \Updownarrow (C \Updownarrow C) \Updownarrow (G \Updownarrow G) = 0$ by commutativity of XOR. This would prove the general case of $\partial^2()$ under this interpretation. \smnl
\tb Now it might be meaningful to look at the points during the decomposition in that for example $CD = C \wedge D$ explicitly. We would then have that $GD = (1,1,0)$, $CD = (0,0,0)$, $CG = (0,0,1)$. Now if $-(CG) = (1,1,1)$ then we have our $(1,1,0) \Updownarrow (1,1,1) \Updownarrow (0,0,1) = (0,0,0)$ \smnl
\tb For completeness $\partial(CGDB) = (GDB) - (CDB) + (CGB) - (CGD)$ Now $(GDB) = (0,1,0)$, $- (CDB) = -(0,0,0) = (1,1,1)$, $ (CGB) = (0,0,0)$, and $ - (CGD) = -(0,0,0) = (1,1,1)$. Now as there is a single $\partial$ and we note that there is not a symmetry as no matter whether $-(CDB)$ and $-(CGD)$ = $(0,0,0)$ or $(1,1,1)$ we have that $(GDB) = (0,1,0)$ So the overall XOR terms for $\partial(CGDB) = (0,1,0)$. So there is an asymmetry in the index that corresponds to $Y$ in $(X,Y,Z)$. Note that this could also be found in $C \Updownarrow G \Updownarrow D \Updownarrow B$. It might be useful to see if the components with negative signs always cancel out no matter if they are $-$ or $+$, thus it wouldn't matter if we take $\Updownarrow$ or $\Leftrightarrow$. \smnl
\tb Now we investigate points and their relationships on an n-dimensional plane! What about then for something such as $D \vee E \vee F \vee G = \mathcal{L}$ so for $\mathcal{V}^T{L}$ we have the points $\{(1,1,0),(1,0,1),(0,1,1),(1,1,1)\}$ associated to the above logical variables in $\mathcal{L}$. Note however that each of the logical points is in terms of $(\mathcal{V}^T(X),\mathcal{V}^T(Y),\mathcal{V}^T(Z))$. Now we might have too many points (for the number of distinct variables) i.e. 4 points for a 3-D object and as each of the points seem to be basis points as no 3 points are co planar i.e. $D\begin{matrix}
1 & 1 & 0 \\ 1 & 0 & 1 \\ 0 & 1 & 1  \end{matrix} = 1(-1) = 1 (1) = -2 \neq 0$ and for $\alpha(1,1,0) + \beta(1,0,1) + \gamma(0,1,1) = (1,1,1) \Rightarrow \alpha + \beta = 1, \alpha + \gamma = 1, \beta + \gamma = 1 \Rightarrow \beta = 1 - \alpha$ so $1 = \alpha + \gamma = 1$ then $- \alpha + \gamma = 0 \Rightarrow \alpha = \gamma$ so $2\gamma = 1 \Rightarrow \gamma = \frac{1}{2}$, $2 \alpha = \frac{1}{2}$ and $\beta = \frac{1}{2}$ so $\alpha,\beta,\gamma =\neq 0$ and we don't need to have all four points to make the simplex. In fact, my algorithm decomposes $X \wedge Y \wedge Z = (1,1,1)$ to $\neg XYZ \Updownarrow X\neg YZ \Updownarrow XY\neg Z$! This is why we can use the 3 non co-planar points (any 3) to give the other! For example with $G = (1,1,1)$, $F = (0,1,1)$, $E = (1,0,1)$, and $D = (1,1,0)$ that for $\{(1,1,1), (0,1,1),(1,0,1)\}$ we have that $F \Updownarrow E = D$. However with 3 points we can calculate each of the other points. So we may have a redundency in points if we take each $n + 1$ in an n-simplex what if then we had 3 co-planar points? say $\{0,1,0), (0,1,1),(0,0,1)\}$? This is then from the logical expression $\neg X \wedge Y \wedge \neg Z \vee \neg  \wedge Y \wedge Z \vee \neg X \wedge \neg Y \vee Z = \neg X \wedge Y \vee \neg X \wedge Z $ which are two points in a 3-D space. Also, what about $\partial(\neg X \wedge Y \wedge Z)$ this is also $= \neg X \wedge Y \vee \neg X \wedge Z $. Now just taking $\neg X \wedge Y \wedge \neg Z \vee \neg X \wedge \neg Y \wedge Z = \neg X \wedge (Y \Updownarrow Z)$ but I shouldn't have said $\vee$ as the points in the set are XOR-ed. So $\partial(0,1,1) = \partial(\neg X \wedge Y \wedge Z) = \neg X \wedge Y \wedge Z \Updownarrow \neg X \wedge \neg Y \wedge Z \Updownarrow \neg X \wedge Y \wedge \neg Z =$ ? Should work. So if we have a logical expression $\neg X \neg Y Z \Updownarrow \neg X Y \neg Z$ this is homologous to $\neg X \wedge (Y \Updownarrow Z)$?? So we can create a "line" of possibilities or perhaps its a plane with points $(0,1,1),(0,0,1)$ and $(0,1,0)$ and if have 3 non co-linear points we can create a plane - or perhaps a 3-simplex tetrahedron of possible positive valuations or "truths".


\end{document}